% Generative Reconciliation Perception
% Wisdom Happy
% 2026

\documentclass[12pt]{article}

\usepackage{amsmath, amssymb, amsthm}
\usepackage{hyperref}
\usepackage{booktabs}
\usepackage{enumitem}
\usepackage[margin=1in]{geometry}

% Theorem environments
\newtheorem{theorem}{Theorem}[section]
\newtheorem{lemma}[theorem]{Lemma}
\newtheorem{corollary}[theorem]{Corollary}
\newtheorem{proposition}[theorem]{Proposition}
\theoremstyle{definition}
\newtheorem{definition}[theorem]{Definition}
\newtheorem{principle}{Principle}
\newtheorem{example}[theorem]{Example}
\theoremstyle{remark}
\newtheorem{remark}[theorem]{Remark}

% GRP notation shortcuts
\newcommand{\grp}{\textsc{grp}}
\newcommand{\imag}{\mathcal{I}}                  % imagined state
\newcommand{\obs}{\mathcal{O}}                   % observed state
\newcommand{\rec}{\mathcal{R}}                   % reconciler output
\newcommand{\err}{\Delta}                        % prediction error

\title{Generative Reconciliation Perception: \\
A Predictive-Imagination Architecture \\
for Computer-Use Cognition}

\author{Wisdom Happy\thanks{Correspondence: wisdomhappy@playfulsincerity.org. \\[0.5em]
\textit{Methodology:} The architecture described in this paper was originally
conceived during a research sprint on March 31, 2026, under the codename
\emph{Phantom}, and developed through April 2026 using the Playful
Sincerity Digital Core---an AI-assisted research methodology system built on
Claude Code (Anthropic). The Digital Core orchestrates structured research,
hierarchical planning, and parallel agent-based exploration. Eight parallel
research agents surveyed approximately 65 academic papers and 80 GitHub
repositories to validate the originality of the architecture's claimed
contributions; results are catalogued in
\texttt{research/landscape.md} and \texttt{research/references.md} in the
project repository. The codename was changed to \grp{} (Generative Reconciliation Perception) on April~22, 2026, to name the actual mechanism rather than a stealth side property,
alongside a sharpening that introduced end-state imagination and the
Generative Collaborative reframe (Section~\ref{sec:gencol}). The complete
specification (\texttt{SPEC.md}, $\sim$640 lines), architecture diagrams,
and project chronicle are publicly available at
\url{https://github.com/Playful-Sincerity/GRP-Generative-Reconciliation-Perception}.
This paper is the first standalone treatment of the architecture; it
complements but does not replace the technical specification.} \\[0.3em]
Playful Sincerity Research}

\date{April 2026}

\begin{document}

\maketitle

\begin{center}
\textit{``The brain is a prediction machine. \\
Perception is controlled hallucination.''}\\[0.3em]
--- standard summary, predictive processing literature
\end{center}
\vspace{1em}

\begin{abstract}
Modern computer-use agents---browser controllers, desktop automators,
screen-grounded language models---share a perceptual posture inherited
from the earliest scripting frameworks: \emph{observe, then interpret}.
The page or screen arrives as input; the model decides what it is and
what to do about it. Imagination, when it appears at all, lives in
chain-of-thought reasoning \emph{between} steps. We propose a different
posture. The agent imagines what it expects to see \emph{before} it
looks. Observation does not produce interpretation directly; it produces
a comparison against the imagined state, and the gap between the
two---the prediction error---is the signal that drives attention,
learning, and downstream action selection. We call this architecture
\textbf{Generative Reconciliation Perception (\grp{})}. The contribution
is not the addition of imagination as a feature. It is the structural
claim that imagination is the substrate on which perception runs.

The architecture supports two modes of imagination. \emph{Next-state
imagination} generates the immediate consequence of the next
action---the WebDreamer-style baseline used to gate when the world
model updates and which tier of perception runs. \emph{End-state
imagination} generates the target final state, optionally as a full
video of the action sequence, and the system reverse-engineers the
concrete actions whose execution would reproduce that target. The
pipeline shifts work to where modern foundation models are
strongest---imagining final states---and away from where they are
weakest---grounding intermediate steps. Two architectural moves complete
the contribution: dual-channel perception (visual and structural) with
a \emph{collaborative} reconciler that treats inter-channel disagreement
as informative cognitive dissonance rather than as error to be
suppressed, and a reframe of the GAN-style adversarial paradigm into a
\emph{Generative Collaborative} one---the imagined world and the
observed world cooperate to converge rather than competing.

This paper specifies the architecture, names the two modes of
imagination, situates \grp{} against the prior work in world models
and predictive processing, and surfaces the open questions whose
resolution will determine which contributions are paper-citable
architecture and which remain useful framing. Implementation has not
started; the architecture is the artifact at this stage. Four
falsification predictions are stated sharply enough to be checked
once implementation begins (Section~\ref{sec:falsification}).
\end{abstract}

% ============================================================
\section{Introduction}
\label{sec:intro}
% ============================================================

\subsection{The Problem}
\label{subsec:problem}

Modern computer-use agents fall into three paradigms. The oldest---Selenium,
Puppeteer, classical scripting---treats the page as a machine to send
commands to. The dominant current paradigm---browser-use, Stagehand, Atlas,
Comet---serializes the page or screen into a textual or visual representation,
prompts a language model with that representation alongside the task,
parses an action out of the response, and executes it. The frontier---UI-TARS,
Lux, OmniParser-driven systems---adds dual-channel perception (visual plus
structural) and improves the grounding pipeline.

What unifies all three is the order of operations. Input arrives first.
The agent decides what it is. The agent acts. The next observation
overwrites the previous one. There is no expectation, no memory of what
was supposed to be there, no internal model that gets surprised. If the
page changed in an unexpected way, the agent has no mechanism for
noticing---only for re-grounding. Each step starts from scratch.

Three structural problems follow.

\begin{enumerate}[leftmargin=*]
\item \textbf{The agent is blind between actions.} Actions take time.
  Pages load, animations run, modals appear, transitions happen. The
  reactive agent has no model of what it expects during this interval, so
  it either polls the screen at high cost or guesses when to move on.
  Time-to-grounding is high; cost is high; reliability is low.
\item \textbf{Surprise is invisible.} When something unexpected
  happens---a layout shift, a new error state, a previously unseen
  element---the agent has no way to register the surprise as such.
  Re-grounding produces a new representation; the \emph{change} between
  old and new representations is not part of the architecture's
  vocabulary. The most informative moments in a session are precisely
  the moments the architecture cannot see.
\item \textbf{Quality scales with chain-of-thought density.} When the
  agent has to reason its way to a correct action, performance on
  complex tasks tracks the depth of the prompt scaffold rather than the
  depth of the model's understanding. This is expensive and brittle.
\end{enumerate}

\subsection{A Precedent: Predictive Processing in Neuroscience}
\label{subsec:precedent}

The structural move \grp{} makes is not new. It is new only as applied
to computer-use agents. Predictive processing---the framework in which
the brain is treated as a generative model whose primary job is to
\emph{predict} sensory input rather than to react to it---has been the
dominant theory of perception in computational neuroscience for two
decades. The free-energy principle (Friston and successors) gives a
formal account in which perception is the minimization of prediction
error between the brain's internal generative model and incoming
sensory input. Action, in the same framework, is the resampling of the
sensory input by moving the body so the prediction is confirmed. The
two are unified: perception and action both reduce free energy.

The empirical case is strong. Visual illusions, top-down attention
effects, repetition suppression in fMRI signals, neural responses to
unexpected stimuli---all are best explained by a brain that imagines
what should be there and reconciles against what is. Animals do not
perceive; they predict and reconcile. The brain is generative, and
observation is the residual the generative model is trying to minimize.

What \grp{} does is apply this architecture to a domain---computer-use
agents---where it is, as of this writing, completely unoccupied in the
published literature. World-model work in browser agents
(WMA~\cite{wma}, WebDreamer~\cite{webdreamer},
MobileDreamer~\cite{mobiledreamer}, CUWM~\cite{cuwm}) does the
next-state half. Active inference for web navigation
specifically~\cite{efe2025} is gestured at but not implemented in any
shipped agent we surveyed. The architectural commitment to imagination
\emph{as substrate}---rather than as one among several reasoning
patterns the model can be prompted into---is, to our knowledge, novel
in the computer-use agent space.

\subsection{\grp{}'s Proposal}
\label{subsec:proposal}

\begin{principle}[Imagination as Substrate]
\label{prin:substrate}
Imagination is not a meta-cognition trick layered on top of perception.
It is the substrate. Perception is the \emph{reconciliation} of imagination
with observation, and prediction error is the primary learning signal
of the system.
\end{principle}

A tighter way to say this: today's computer-use agents perceive in order
to act. \grp{} imagines in order to perceive, and perceives in order to
confirm or correct what was imagined. Imagination is not the optional
flourish a sufficiently clever model adds when it has spare reasoning
to spend. It is the layer the architecture runs on. Every observation
is a comparison; every comparison produces a residual; every residual
is a learning signal.

Three claims follow from the principle.

\textbf{First, imagination is best treated as infrastructure, not as
inference.} When imagination is implemented as chain-of-thought
reasoning during prompting, it competes with action selection, error
recovery, and other reasoning the model is trying to do in the same
context. When imagination is a structural component---a dedicated
engine with explicit inputs (memory, task, last action) and explicit
outputs (expected state, end-state target, confidence)---it stops
competing with reasoning and starts subsidizing it. The agent reasons
less about \emph{what should be there} because the imagination engine
answers that question before reasoning begins.

\textbf{Second, the high-value use of imagination is end-state, not
next-state.} Most prior work on predictive perception in agents
generates the immediate next state. This is useful: it gates
learning, cheapens reconciliation, enables tiered perception. But the
more valuable mode imagines the \emph{final} state---the target page
after the design change has been applied, the final video frame after
the robotic task is complete---and reverse-engineers the action
sequence whose execution would reproduce that target. This inverts
the standard pipeline. Quality of output scales with quality of
imagination; the model's reasoning load shifts from \emph{grounding
intermediate steps} (where it is brittle) to \emph{imagining final
states} (where it is strong).

\textbf{Third, reconciliation between imagination and observation is
collaborative, not adversarial.} GAN-style architectures frame
generation and discrimination as competition: one tries to fool, one
tries to detect. \grp{} frames the imagined world and the observed
world as collaborators working to converge. The agent adjusts its
actions to bring the observation toward the imagined target, and
adjusts its imagination toward what it actually observes when surprise
indicates the imagination was wrong. Both sides update. Surprise is
information, not failure.

\subsection{Summary of Results}
\label{subsec:summary}

This paper contributes:

\begin{enumerate}[leftmargin=*]
\item A specification of imagination-as-substrate as an architectural
  commitment for computer-use agents (Section~\ref{sec:loop}), not as
  a chain-of-thought pattern.
\item Two modes of imagination, with end-state imagination plus
  reverse-engineering identified as the load-bearing
  contribution (Section~\ref{sec:two-modes}).
\item A dual-channel perception architecture with a collaborative
  reconciler that treats inter-channel disagreement as first-class
  informative signal (Section~\ref{sec:dual}). No published browser-agent
  framework treats disagreement this way.
\item The Generative Collaborative reframe of the GAN-style adversarial
  paradigm, with explicit acknowledgment that a formal loss-function
  statement is open work (Section~\ref{sec:gencol}, Q4).
\item A tiered-perception cost-management structure in which imagination
  subsidizes attention and only surprise escalates expensive
  perception (Section~\ref{sec:cost}).
\item An integration story locating \grp{} as the perception subsystem of
  a broader research program (the Synthetic Sentiences Project), where
  the imagination engine is shared infrastructure across perception,
  value-aligned modulation, and dream cycles
  (Section~\ref{sec:integration}).
\item Four sharply stated falsification predictions
  (Section~\ref{sec:falsification}) and five paper-relevant open
  questions (Section~\ref{sec:open}).
\end{enumerate}

% ============================================================
\section{The Cognitive Loop}
\label{sec:loop}
% ============================================================

The full \grp{} loop is:
\begin{equation*}
\begin{aligned}
&\text{IMAGINE} \to \text{OBSERVE} \to \text{RECONCILE} \to \text{DIFF} \\
&\quad \to \text{PLAN} \to \text{CHECK} \to \text{GUARD} \to \text{ACT} \to \text{LEARN}
\end{aligned}
\end{equation*}
with a closure arrow from action back to imagination through the memory
graph. Each stage has a specific functional role.

\begin{description}[leftmargin=*]
\item[\textbf{IMAGINE}] generates either an \texttt{ImaginedState}
  (next-state mode: the immediate consequence of the next action) or an
  \texttt{ImaginedEndState} (end-state mode: the target final state,
  optionally as a video). The engine reads from the memory graph, the
  task context, and the last action.
\item[\textbf{OBSERVE}] runs dual perception: a visual channel
  (screenshot plus Set-of-Mark numbering) and an analytical channel
  (DOM, accessibility tree, macOS Accessibility API), in parallel.
\item[\textbf{RECONCILE}] merges the two channels into a unified
  \texttt{PageModel}. Channel disagreements are surfaced as cognitive
  dissonance---first-class structural information---and resolved by
  collaborative arbitration rather than by suppressing one channel
  in favor of the other (Section~\ref{sec:dual}).
\item[\textbf{DIFF}] computes prediction error
  $\err = d(\imag, \obs)$: how far is the imagined state from the
  observed \texttt{PageModel}? In end-state mode, the diff also
  computes how far the observed state is from the target.
\item[\textbf{PLAN}] selects the next action from an affordance-aware
  view of the page. Elements are perceived as action potentials
  (Gibson's ecological psychology) rather than as inert structure.
  Spreading activation through the memory graph weights affordances
  by relevance to the task.
\item[\textbf{CHECK}] runs the alignment critic---a cheap Haiku-class
  call evaluating the candidate action against the original task
  intent. The critic exists to interrupt the loop when the agent's
  action drifts from what was asked, not because the model has gone
  rogue but because the prediction-error path is locally optimal in
  ways the task may not endorse.
\item[\textbf{GUARD}] enforces the security membrane: content
  isolation, HTTP-layer sandboxing, credential isolation, permission
  tiers.
\item[\textbf{ACT}] executes the action through the motor system, with
  human-like input dynamics where the environment requires it.
\item[\textbf{LEARN}] writes back to the memory graph if the prediction
  error exceeds a threshold. Most observations do not produce a write.
  Memory consolidates surprise, not routine.
\end{description}

The closure arrow---from action back to imagination---is the cycle
that makes the system \emph{learn from experience}. Every action
changes the world; every observation of that change produces a
residual against the imagined consequence; every residual that
exceeds threshold updates the model that generated the imagination
in the first place. The agent's world model is the cumulative
residual of its own predictions over time. Reconsolidation is the
mechanism by which it improves.

% ============================================================
\section{The Structure of Imagination}
\label{sec:two-modes}
% ============================================================

Before describing the two modes, it matters to be specific about what
the imagination engine \emph{is} computationally. Imagination, in
\grp{}, is not a metaphor for chain-of-thought. It is a typed
generative procedure with explicit inputs and outputs, implemented by
specific model classes appropriate to the granularity of what is being
imagined.

\subsection{What the Imagination Engine Generates}
\label{subsec:structure-of-imagination}

The imagination engine is polymorphic over output type. Different
imagination targets call for different generative substrates.

\paragraph{Structural imagination.}
For next-state prediction (Section~\ref{subsec:next-state}) and for
abstract end-state targets where layout and structure matter more than
pixels, the engine queries a foundation language model conditioned on
the memory graph, the task context, and the last action. Output is a
typed \texttt{ImaginedState}---expected URL pattern, expected elements
with positions and affordances, layout expectations, confidence with
provenance. This is the WebDreamer-style baseline. Cheap, fast, and
sufficient when the question is ``what should be there?'' rather than
``exactly what should it look like?''

\paragraph{Visual imagination via diffusion models.}
For end-state targets where pixel-accurate visual fidelity is the
specification---a target page mock, a target frame in a robotic task,
a finished design---the engine generates the imagined target using
image or video diffusion models. The agent imagines a screenshot
literally. It produces an image or a sequence of frames. This image
\emph{is} the target the reverse-engineering pass closes against.
The choice of substrate matters: diffusion models are trained on
exactly the corpus of finished visual artifacts the agent needs to
imagine, and they generate at the resolution the dual-perception
visual cortex consumes. The same model class that generates art is
the model class that generates the target the agent is trying to
build.

\paragraph{Temporal imagination.}
For tasks where the target is a process rather than a state---a video
of the user experiencing the page correctly, the trajectory of a
robotic motion---the engine generates frame sequences via video
diffusion or autoregressive video models. Reverse-engineering then
applies inverse dynamics (for motor) or differential codegen (for
visual UI) to derive the action sequence whose execution reproduces
the imagined frames.

\bigskip
\noindent
The architectural commitment is that imagination is a \emph{component}
with multiple implementations, not a single inference pass. The same
loop runs over all three substrate choices; what changes is the
specific generative model invoked for the imagine step. This factoring
is what lets the architecture absorb improvements in diffusion models,
video generation, and language-as-world-model independently. As any
of these substrates improves, \grp{}'s imagination quality improves;
the rest of the loop does not need to change.

\subsection{Next-State Imagination}
\label{subsec:next-state}

Given the current state of the page or screen, the most recent
action, and the agent's memory of similar contexts, generate the
expected next state. This is the WebDreamer-style baseline. It runs
on every action and serves three purposes: it gates the perception
tier (Section~\ref{sec:cost}); it gates the memory write (if the
prediction was correct, no reconsolidation is needed; if it was
wrong, the surprise is encoded); and it cheapens reconciliation (the
reconciler has a strong prior to compare against, and inter-channel
disagreement is easier to interpret when the expected state is
known).

The contribution at this layer is not novelty in the imagination
itself---WMA, WebDreamer, MobileDreamer, and CUWM all do versions of
this. The contribution is treating imagination as
\emph{infrastructure} rather than as inference: it runs always, it
produces a typed object (\texttt{ImaginedState}), and downstream
components depend on it as a structural input. Cost is amortized
across the whole loop.

\subsection{End-State Imagination and Reverse-Engineering}
\label{subsec:end-state}

Given the task intent and the current state, imagine the
\emph{target}---the final state of the page or screen after the task
is complete, or, more powerfully, an entire video of the action
sequence playing out correctly. Then reverse-engineer the action
sequence whose execution would reproduce that target.

This is the load-bearing contribution of the paper. Two concrete
instantiations make the move clear.

\paragraph{Web design.}
The current page is a screenshot. The task is ``make this page match
the design intent.'' The imagination engine generates a target
screenshot---what the page should look like after the change---or a
short video of a user experiencing the page as it is supposed to be
experienced. Reverse-engineering takes that target and produces the
diff in HTML, CSS, and JavaScript that closes the gap. The agent
runs the diff, observes the result, reconciles it against the
imagined target, and iterates. The pipeline shifts work to where
modern foundation models are strongest---imagining final
states---and away from where they are weakest---grounding
intermediate code edits.

\paragraph{Robotics and motor control.}
The task is ``pick up the cup and pour it.'' The imagination engine
generates a video of the task being performed correctly: the
trajectory of the arm, the grip on the handle, the angle of the pour,
the rest. Inverse dynamics derives the motor commands whose execution
reproduces those frames. The agent executes the commands, observes
the actual video, reconciles it against the imagined video, and
updates either the motor model (if the actions did not produce the
imagined effects) or the imagination engine itself (if the imagined
target was off-distribution from what the embodiment can actually
do).

The mechanism is the same in both cases. The output channel---code or
motor commands---is different. The pipeline structure is
\emph{imagine target $\to$ derive actions $\to$ execute $\to$
reconcile against target $\to$ iterate}. This is not a chain-of-thought
trick. It is an architectural inversion of how the agent decides what
to do. Quality of output scales with quality of imagination, because
the imagination \emph{is} the specification.

\subsection{Why End-State Matters More}
\label{subsec:why-end-state}

End-state mode plays to the model's strengths: foundation models---and
diffusion models in particular---are trained on enormous corpora of
finished artifacts, and they imagine final states with detail and
consistency. They are much weaker at intermediate-step reasoning, where
errors compound and brittleness scales with depth. End-state imagination
skips the brittle middle. It also compresses long-horizon planning:
fifty actions are hard to plan as fifty decisions but easy to plan as
one imagined target plus an inverse procedure. And it produces
verifiable artifacts---the imagined target is something the agent can
\emph{show}, making intent legible to a human reviewer in a way
intermediate reasoning chains usually are not. The trade-off is
computational: imagining a full thirty-second video at frame rate is
expensive, and sparse keyframes, structural targets, and progressive
refinement are open optimizations (Section~\ref{sec:open}, Q1).

% ============================================================
\section{Dual Perception with Collaborative Reconciliation}
\label{sec:dual}
% ============================================================

\grp{} runs two perception channels in parallel.

\subsection{Visual Cortex}
A screenshot of the page or screen is processed by a vision-language
model, with Set-of-Mark numbering applied to interactable elements.
The output describes the spatial layout, the visual states of
elements (highlighted, disabled, loading, error), the colors and
typography, the overall composition. \emph{Strengths:} sees what is
actually rendered, including custom widgets, dynamic content, and
visual feedback that is not exposed in structured form.
\emph{Weaknesses:} grounding is approximate, exact text values are
sometimes hallucinated, and the channel is expensive.

\subsection{Analytical Cortex}
The DOM, accessibility tree, or macOS Accessibility API is queried
directly. The output describes the structural layout: element roles,
ARIA labels, exact text values, click handlers, form state.
\emph{Strengths:} ground truth on structure and interactivity, free
or near-free in token cost, deterministic. \emph{Weaknesses:} cannot
see dynamic visual state, custom-painted canvases, animations, or
anything not exposed through accessibility APIs.

\subsection{Disagreement as Information}
The two channels disagree often. The visual cortex sees a button
labeled ``Submit''; the analytical cortex reports a
\texttt{<div role="button">} with no \texttt{aria-label}. The visual
cortex sees a loading spinner; the analytical cortex reports the page
as fully loaded. The visual cortex sees a modal blocking the page;
the analytical cortex reports the modal as hidden behind a
\texttt{display:none} rule that is being overridden by an inline
style. \emph{These are not channel errors. They are information.}
The disagreement itself is the signal.

The reconciler treats inter-channel disagreement as cognitive
dissonance: a structural property of the joint representation, not an
error to be suppressed by picking one channel. Each disagreement is
recorded as a \texttt{Disagreement} object with the visual claim, the
analytical claim, and a resolution rule. The arbitration is
collaborative: visual wins for ``looks right?'' questions, analytical
wins for ``is clickable?'' questions, both are flagged as uncertain
when neither wins, and the disagreement itself becomes a node in the
agent's memory graph for future reference. The next time a similar
disagreement arises on the same site, the agent does not have to
re-derive the resolution.

No published browser-agent framework we surveyed treats inter-channel
disagreement as a first-class informative signal. UI-TARS combines
vision and structure but picks one when they disagree. browser-use
does the same. The published work assumes the channels should agree,
that disagreement is a defect, and that the architecture should hide
it. \grp{} makes the opposite assumption: disagreement is the most
informative thing the dual perception produces, and it should be
surfaced.

% ============================================================
\section{The Generative Collaborative Reframe}
\label{sec:gencol}
% ============================================================

The standard generative-discriminative pattern in machine learning is
adversarial. A generator produces candidates; a discriminator tries to
detect them; the system reaches equilibrium at the saddle point of a
minimax loss. This pattern, formalized in the GAN
literature~\cite{gan2014} and elaborated extensively since, has shaped
how researchers think about generation-plus-evaluation systems for
over a decade.

\grp{} rejects the framing for its own architecture. The imagined
world and the observed world are not adversaries. They are
collaborators working to converge. The agent adjusts its actions so
the observed world bends toward the imagined one; the agent adjusts
its imagination so the imagined world tracks what is actually true
about the observed one. Both sides update. Both sides cooperate. The
reconciler is not a discriminator picking a winner---it is a
collaborative procedure that closes the gap by updating whichever
side is most-cheaply-updated given current evidence.

We call the resulting framing \emph{Generative Collaborative}. The
natural acronym, GCN, is taken in the machine-learning literature
(Graph Convolutional Networks); a paper publication will need a
different shorthand, but the concept stands regardless of how it is
eventually labeled.

\begin{remark}[Status of the reframe]
\label{rem:gencol-status}
The Generative Collaborative reframe is currently a \emph{framing
decision} backed by an architectural commitment, not a citable
architectural contribution with a formal loss-function statement.
Until the formal mathematical content of ``cooperate''---single shared
loss, two losses provably aligned in expectation, or something more
complex---is grounded, the reframe should be presented as the framing
under which \grp{}'s reconciler is designed, not as a published
mathematical result. This is the most important open question in the
paper (Section~\ref{sec:open}, Q4).
\end{remark}

Three things follow from the reframe.

\textbf{(1) Surprise is not failure.} In a GAN, the generator failing
to fool the discriminator means the generator must improve. In \grp{},
the imagined state failing to match the observed state means the world
model has new information. Surprise is the input the reconciler is
\emph{for}. Suppressing surprise---picking a winner channel, hiding
disagreements, training the imagination to never be wrong---would
defeat the architecture.

\textbf{(2) The two sides share a loss.} What loss, exactly, the
imagination and the observation share is open
(Remark~\ref{rem:gencol-status}). The framing says they cooperate; the
formal mathematical content of ``cooperate'' is not yet settled.

\textbf{(3) The framing is consistent with a broader research
program.} \grp{} is one subsystem of a larger project (the Synthetic
Sentiences Project) whose alignment thesis is that aligned behavior
emerges from architectures whose internal substrates \emph{converge
through attraction} rather than being constrained by external
penalties. The collaborative reconciler is the same move at the
perception layer that the broader project makes at the
value-alignment layer: cooperation rather than competition. This is
not coincidence; it is design.

% ============================================================
\section{Cost Management: Tiered Perception}
\label{sec:cost}
% ============================================================

The full \grp{} pipeline is expensive: imagination, dual perception,
reconciliation, alignment critic, motor system, memory write-back.
Running it on every action would defeat the architecture for any task
with more than a handful of steps.

The architecture defines four tiers of perception, selected by the
prediction error from the previous step (Table~\ref{tab:tiers}).

\begin{table}[h]
\centering
\begin{tabular}{lllll}
\toprule
\textbf{Tier} & \textbf{Name} & \textbf{What runs} & \textbf{Tokens} & \textbf{When} \\
\midrule
0 & Reflex & DOM-only check, no LLM call & 0 & Confirming, waiting \\
1 & Glance & DOM + light a11y, Haiku-class & $\sim 500$ & Familiar pages \\
2 & Look & Screenshot + SoM + a11y, Sonnet & $\sim 2{,}000$ & New pages \\
3 & Study & Full pipeline + imagination & $\sim 5{,}000$ & First visit, high surprise \\
\bottomrule
\end{tabular}
\caption{Tiered perception. Tier escalates on prediction error.}
\label{tab:tiers}
\end{table}

A ten-step task with mostly familiar transitions runs at Tier~0 or~1
for most steps and at Tier~3 only at the moments the agent is
genuinely uncertain. The cost of a typical task drops from
$\sim 50{,}000$ tokens (running the full pipeline every step) to
$\sim 15{,}000$ tokens. The architecture is not just better; it is also
cheaper to run, because imagination subsidizes attention.

This tier structure is the same Kahneman dual-process pattern applied
to perception. Familiar inputs run on System~1 (cheap, automatic).
Surprising inputs escalate to System~2 (expensive, deliberate).
Imagination is what makes the distinction possible---without an
expectation, every observation looks equally novel, and there is
nothing to escalate from.

% ============================================================
\section{Architecture in Brief}
\label{sec:architecture}
% ============================================================

The full architectural specification is in \texttt{SPEC.md}
($\sim$640 lines, including data structures, security model, and
module map). What follows is the structural summary at the level
needed for this paper.

\subsection{Data Structures}

\texttt{ImaginedState} (next-state mode) carries the expected URL
pattern, expected elements with positions and affordances, layout
expectations, and a confidence value with provenance.

\texttt{ImaginedEndState} (end-state mode) carries a target screenshot
or video sequence, the intent, a confidence value, and a
reverse-engineering strategy: visual-diff-then-codegen, motor-inverse,
or plan-search.

\texttt{PageModel} is the reconciled output of dual perception:
elements with confidence levels, layout, disagreements, and
affordances.

\texttt{Affordance} describes an element's action potential---type,
expected outcome, confidence, relevance to task, risk class, source.

\texttt{Disagreement} describes inter-channel conflicts with their
resolution rule and provenance.

\subsection{Memory}

The world model is not part of \grp{}. It lives in a separate
subsystem (\textsc{mwm}: Memory as World Model), and \grp{} writes
content into \textsc{mwm}'s primitives rather than owning a graph
schema. This factoring is deliberate: every subsystem in the broader
project that produces or consumes beliefs touches the same graph, and
\grp{}'s perception flood is one input among several
(Section~\ref{sec:integration}).

\subsection{Pluggable Backends}

Perception is pluggable across Claude Vision, Lux (OpenAGI),
OmniParser, Playwright accessibility, and macOS Accessibility API.
Action is pluggable across Playwright, Patchright, Camoufox,
nodriver, PyAutoGUI, and macOS AppleScript. Adapters cover browser,
desktop, and any-window environments.

\subsection{Security Membrane}

Filesystem sandbox at \texttt{\textasciitilde/.grp/workspace/};
content isolation between web text and agent instructions; HTTP-layer
sandbox with domain allowlist and method restrictions; credential
broker that fills directly through the browser engine without exposing
credentials to the LLM context; alignment critic on every action;
permission tiers (auto, confirm, deny) governing different action
classes.

\subsection{Build Phases}

Implementation proceeds in six phases:
\emph{Walking} (basic Playwright + visual perception),
\emph{Seeing Double} (dual perception + reconciler),
\emph{Dreaming} (next-state imagination),
\emph{Imagining the End} (end-state imagination + reverse-engineering),
\emph{Feeling} (temporal perception + motor system),
\emph{Speaking} (alignment critic + full security membrane).

End-state imagination is intentionally introduced \emph{after} the
next-state loop is working, because reverse-engineering depends on a
stable substrate to verify against.

% ============================================================
\section{Where \grp{} Belongs}
\label{sec:integration}
% ============================================================

\grp{} does not stand alone, even though this paper presents it as a
standalone contribution. It is the perception subsystem of a broader
research program---the Synthetic Sentiences Project---whose other
subsystems consume \grp{}'s outputs and shape what \grp{} imagines.

\paragraph{Memory.}
The world model lives in \textsc{mwm}~\cite{mwm}. \grp{} writes
predicted states, observed states, and causal edges (clicking $X$ led
to state $Y$) into \textsc{mwm} as nodes and edges. The graph is the
agent's long-term substrate; \grp{} is the perceptual input layer that
feeds it. When the agent asks ``have I seen this site before?'' or
``what usually happens after this action?'', the answer is a
traversal of the \textsc{mwm} graph, not an internal-to-\grp{} lookup.

\paragraph{Earned conviction.}
Page elements have confidence levels rooted in evidence chains. The
dissonance-as-information primitive in \grp{}'s reconciler is the same
primitive the cognition subsystem uses to surface unresolved beliefs.
Curiosity about unknown elements (low-confidence nodes) is the same
affect that drives investigation in the broader cognitive layer.

\paragraph{Value-aligned modulation.}
What to attend to, what to imagine, and what counts as an acceptable
action are weighted by the agent's current value-alignment state. The
alignment critic that evaluates each action before execution is a
value-check at the perception/action boundary. The imagination engine
itself is conditioned on what the agent cares about: the same task
can produce different imagined end-states depending on the values
active.

\paragraph{Working memory and dream cycles.}
The page being perceived lives in a working tree that grows when new
perceptual content is loaded and shrinks when relevance fades. Sleep
cycles consolidate \grp{}'s perceptual outputs into long-term graph
structure. Dream cycles use \grp{}'s imagination engine to
\emph{simulate} without observing---adversarial self-reflection,
value-conflict tests, hypothetical scenarios. Imagination is the
substrate of perception when the agent is awake and the substrate of
simulation when the agent is asleep.

\paragraph{Voice.}
When the agent speaks about what it perceived, the acoustic prosody
of its voice should reflect the confidence levels in the page model.
Honest disfluency around low-confidence perceptions is prosodic
accuracy, not a speech defect.

\bigskip
\noindent
The factoring is deliberate. The imagination engine \grp{} uses for
end-state targets, the should-world model that drives value-aligned
modulation, and the simulation substrate dream cycles run on are the
\emph{same mechanism} applied at three scales: task-scale (\grp{}),
being-scale (values), moment-scale (cycles). One imagination engine,
three uses. This is the unification thesis the broader project rests
on.

% ============================================================
\section{Literature Positioning}
\label{sec:literature}
% ============================================================

\grp{} draws from and stands against four overlapping bodies of work.

\subsection{World Models for Computer-Use Agents}
WMA~\cite{wma} introduced transition-focused observation abstraction
for predicting how actions change pages.
WebDreamer~\cite{webdreamer} used the LLM itself as a world model,
``dreaming'' outcomes before committing---closest in spirit to
\grp{}'s next-state mode but with no end-state mode, no reconciler,
and no architectural integration with memory.
MobileDreamer~\cite{mobiledreamer} showed sparse text sketches
suffice for effective world modeling.
WebWorld~\cite{webworld} and CUWM~\cite{cuwm} advanced trajectory
simulation and two-stage prediction respectively.
DreamerV3~\cite{dreamerv3} provides the foundational world model
architecture from RL.

\grp{}'s contribution against this line is end-state imagination with
reverse-engineering as a unified architectural mode, not an emergent
capability.

\subsection{Visual Grounding for Agents}
SeeAct~\cite{seeact} established that grounding is the bottleneck,
not reasoning.
Set-of-Mark~\cite{som} and OmniParser~\cite{omniparser} provide
standard tools for converting screenshots into addressable elements.
UGround~\cite{uground} showed universal visual grounding generalizes
across UIs.
WebVoyager~\cite{webvoyager} demonstrated a 28-point gap between
SoM-based and text-only agents.

\grp{} uses these tools as the visual channel of dual perception; the
contribution is the reconciler that combines them with structural
input rather than picking between them.

\subsection{Active Inference and Predictive Processing}
The free-energy principle and predictive processing literature
(Friston and successors) provide the theoretical grounding for
imagination-then-observation as a perceptual architecture.
\cite{efe2025} shows expected free energy reduces to prediction error
in the relevant limit.
\cite{causal2024} proves causal world models are necessary for
robust generalization.
\cite{causalang2024} connects causal variables to language.

\grp{} applies this framing to a domain---browser and computer-use
agents---where it is, as of this writing, completely unoccupied in
the published literature.

\subsection{Generative Adversarial Architectures}
The GAN literature~\cite{gan2014} establishes the dominant pattern
for generation-plus-evaluation systems. \grp{} rejects the framing
for its perception layer (Section~\ref{sec:gencol}). The Generative
Collaborative reframe is, to our knowledge, novel as an architectural
application to perception---though related ideas exist in cooperative
multi-agent learning.

% ============================================================
\section{Falsification}
\label{sec:falsification}
% ============================================================

The architecture commits to specific predictions. Four are sharp
enough to falsify.

\begin{description}[leftmargin=*]
\item[\textbf{F1.}] \emph{End-state imagination outperforms
  chain-of-thought for tasks of complexity above a threshold.} On a
  benchmark of web-design and motor-control tasks of varying
  complexity (measured by step count and error compounding), end-state
  imagination plus reverse-engineering should produce higher success
  rates than chain-of-thought reasoning at matched compute, with the
  gap growing with complexity. \emph{If end-state imagination
  underperforms or matches chain-of-thought across the complexity
  spectrum, the central claim of the paper is wrong.}

\item[\textbf{F2.}] \emph{Disagreement between dual perception
  channels is informative.} Channel disagreements should correlate
  with locations in the page where prediction error is later highest,
  with elements that subsequently produce action errors, and with
  sites where causal-edge updates in the memory graph are
  highest-weighted. \emph{If disagreements are random noise,
  uncorrelated with downstream signal, the reconciler's
  dissonance-as-information claim is empirically unsupported.}

\item[\textbf{F3.}] \emph{Imagination subsidizes perception cost.} A
  complete \grp{} pipeline running on a benchmark of mixed-difficulty
  tasks should produce a token-cost distribution where the median
  step runs at Tier~0 or~1, escalating to Tier~3 only at moments of
  high prediction error. \emph{If the cost distribution is flat across
  tiers, the architecture is not actually subsidizing cheap perception
  with imagination---it is just adding cost.}

\item[\textbf{F4.}] \emph{The collaborative reconciler is more
  accurate than a winner-takes-all reconciler.} A controlled
  comparison where one variant of the system uses \grp{}'s
  collaborative arbitration and the other variant picks the
  higher-confidence channel should show the collaborative variant
  produce more accurate \texttt{PageModel} outputs and fewer
  downstream action errors. \emph{If the winner-takes-all variant
  matches or beats the collaborative variant, the dissonance-as-information
  claim is empirically unsupported.}
\end{description}

These predictions cannot all be checked at once; they require an
implementation, a benchmark, and a controlled comparison. They are
stated here so future work can confirm, refine, or reject each.

% ============================================================
\section{Open Questions}
\label{sec:open}
% ============================================================

The architecture is specified. Implementation has not started. Five
questions are paper-relevant; nine are tracked in
\texttt{research/open-questions.md} in the project repository. The
five most important:

\begin{description}[leftmargin=*]
\item[\textbf{Q1.}] \emph{How does end-state imagination scale in
  compute?} A thirty-second video target is hundreds of frames. The
  cost-benefit curve relative to extended chain-of-thought reasoning
  is unknown. Open: the minimum-viable representation of an end-state
  target (single frame, sparse keyframes, structural target plus
  visual hint), and the task complexity at which end-state imagination
  beats chain-of-thought.

\item[\textbf{Q2.}] \emph{How does reconciliation arbitrate when both
  sides are uncertain?} The standard reconciler assumes the observation
  is the more reliable anchor. End-state mode breaks this---the
  imagination \emph{is} the target; observation is what we are trying
  to bend toward it. Both sides have confidence levels that may be
  similar. The right confidence-arbitration rule when the imagined and
  observed sides have comparable confidence is not yet specified.

\item[\textbf{Q3.}] \emph{How does \grp{} write back into memory
  without polluting the graph?} \grp{} produces a flood of perception
  nodes and causal-edge candidates. Not every perception belongs in
  long-term memory. The consolidation threshold, the distinction
  between stable site patterns and one-off page states, and the
  coordination with sleep cycles for staged consolidation are all
  open.

\item[\textbf{Q4.}] \emph{What is the formal mathematical content of
  ``collaborative'' versus ``adversarial''?} The Generative
  Collaborative reframe is currently framing, not architecture-with-a-citable-loss-function.
  Until there is a single shared loss, or two losses provably aligned
  in expectation, the reframe is paper-citable as a framing decision
  but not as an architectural contribution. \emph{This is the most
  important open question in the paper.}

\item[\textbf{Q5.}] \emph{Reverse-engineering reliability across the
  two paths.} End-state to action sequence is two distinct problems
  with different failure modes. The web-design path (visual diff to
  code) is well-studied as image-to-code, but iterative refinement
  under prediction error is open. The motor path (imagined video to
  motor commands) is well-studied as inverse dynamics, but with
  imagined frames as legitimate input rather than observed frames;
  whether the inverse model generalizes when the imagined video is
  slightly off-distribution is an empirical question.
\end{description}

% ============================================================
\section{Limitations and Scope}
\label{sec:limitations}
% ============================================================

\subsection{What \grp{} Does Not Do}

\textbf{It does not eliminate language-model dependence.} The
imagination engine, the reconciler, the alignment critic, and the
planner all rely on foundation models. \grp{} restructures how those
models are used, but does not replace them.

\textbf{It does not solve the grounding problem.} Visual grounding
remains hard. \grp{} mitigates it by adding a structural channel and
a collaborative reconciler, but neither channel is reliable in
isolation, and the reconciler depends on both being approximately
correct most of the time.

\textbf{It does not eliminate the need for benchmarks.} The
falsification predictions in Section~\ref{sec:falsification} require
a controlled comparison against existing systems. The architecture is
a commitment until those comparisons are run.

\subsection{Scope Beyond Computer-Use}

The same architecture---generate expectation, observe, reconcile,
learn from residual---applies beyond browsers: to robotics
(Section~\ref{subsec:end-state}), to scientific instrumentation (the
experiment is the act, the readout is the observation, the prior
model is the imagination), to human-AI collaboration (the agent
imagines what the human wants, observes the reaction, reconciles).
The computer-use case is the most immediately tractable and the
closest to deployment. If the move works there, the architectural
commitment generalizes.

% ============================================================
\section{Outlook}
\label{sec:outlook}
% ============================================================

\grp{} is an architectural commitment, not a finished system. The
specification is complete; the implementation is in front of us. The
most immediate work is the next-state pipeline (Phases 1--3 in the
build plan), the integration with the memory subsystem, and the
first benchmark of imagination-subsidized perception on a
representative web-control task. The end-state mode (Phase 4) follows
the next-state mode is working.

The paper-blocking work is Q4---the formal mathematical content of
the collaborative reframe. Until that is grounded, the Generative
Collaborative claim should be presented as a framing decision and a
research direction, not as a citable architectural contribution. This
is the honest version of where the paper stands.

What we are committing to, in writing: imagination is not a feature
of sufficiently sophisticated reasoning. It is the substrate.
Perception is reconciliation. The architecture's job is to imagine
well, observe carefully, and learn from the gap.

% ============================================================
\section*{Acknowledgments}
\label{sec:acknowledgments}
% ============================================================

The original architecture was conceived under the codename
\emph{Phantom} during a research sprint on March 31, 2026, and was
developed through April 2026 with the assistance of the Playful
Sincerity Digital Core---a configuration of Claude Code rules,
skills, hooks, and chronicles maintained as a research and methodology
partner. The codename was changed to \grp{} (Generative Reconciliation
Perception) on April~22, 2026, to name the actual mechanism rather
than a stealth side property, alongside a sharpening that introduced
end-state imagination and the Generative Collaborative reframe. Eight
parallel research agents surveyed approximately 65 papers and 80
GitHub repositories to validate the originality of the architecture's
claimed contributions; results are catalogued in
\texttt{research/landscape.md} and \texttt{research/references.md}.
The Synthetic Sentiences Project, of which \grp{} is the perception
subsystem, provides the broader architectural context within which
\grp{}'s contributions cohere; the sister project Memory as World
Model (MWM)~\cite{mwm} provides the graph substrate \grp{} writes into.

% ============================================================
\begin{thebibliography}{99}
% ============================================================

\bibitem{wma}
WMA: Web Agents with World Models. arXiv:2410.13232 (2024).

\bibitem{webdreamer}
WebDreamer: LLM as World Model for Web Agents. arXiv:2411.06559 (2024).

\bibitem{mobiledreamer}
MobileDreamer: Textual Sketch World Models. arXiv:2601.04035 (2026).

\bibitem{webworld}
WebWorld: 1M+ Trajectory Simulator for Web Agents. arXiv:2602.14721 (2026).

\bibitem{cuwm}
CUWM: Computer-Using World Model. arXiv:2602.17365 (2026).

\bibitem{dreamerv3}
DreamerV3. arXiv:2301.04104 (2023). Nature (2025).

\bibitem{seeact}
SeeAct: GPT-4V is a Generalist Web Agent. ICML 2024 / arXiv:2401.01614.

\bibitem{som}
Set-of-Mark Prompting Unleashes Extraordinary Visual Grounding in GPT-4V.
Microsoft Research (2023).

\bibitem{omniparser}
OmniParser for Pure Vision Based GUI Agent. arXiv:2408.00203 (2024).

\bibitem{uground}
UGround: Universal Visual Grounding for GUI Agents. ICLR 2025 Oral /
arXiv:2410.05243.

\bibitem{webvoyager}
WebVoyager: End-to-End Web Agent with Large Multimodal Models. ACL 2024 /
arXiv:2401.13919.

\bibitem{efe2025}
Expected Free Energy as Variational Inference. arXiv:2504.14898 (2025).

\bibitem{causal2024}
Robust Agents Learn Causal World Models. arXiv:2402.10877 (2024).

\bibitem{causalang2024}
Language Agents Meet Causality. arXiv:2410.19923 (2024).

\bibitem{gan2014}
Goodfellow, I.~et~al. Generative Adversarial Nets. NeurIPS 2014.

\bibitem{mwm}
Memory as World Model (MWM). Playful Sincerity Research, April 2026. \\
\url{https://github.com/Playful-Sincerity/MWM-Memory-as-World-Model}.

\end{thebibliography}

\end{document}
